% ===================================================================
%  HOTO Na 8 — Girbi. --- Farauta, girbin inabi, kamun kifi.
%
%  Traduction, faite en 2026, en haoussa d'aujourd'hui, sur l'ido de
%  texto/io. Regles de la colonne : voir 10-hoto-01.tex. Deux scenes,
%  deux numerotations.
%
%  ECART DECLARE, ET IL FAUT LE SAVOIR AVANT D'ECRIRE. Le bloc
%  t08-c2-03-1 est l'un des trois que renvoji.py laisse passer sans
%  rien dire : le fac-simile ido y compose « \textsuperscript{13)} »,
%  malforme, sans la parenthese ouvrante, et la machine ne voit donc
%  pas ce renvoi. Un ecart declare rend le bloc AVEUGLE — le controle
%  n'y regarde plus rien du tout. L'ordre juste, pris sur le volume
%  francais, est 3, 12, 14, 13, 15, 16 ; il a ete pose a la main et
%  recompte a la main.
%
%  LES QUATRE POINTS CARDINAUX SONT HAOUSSA, TOUS LES QUATRE : arewa
%  le nord, kudu le sud, gabas l'est, yamma l'ouest. Le vent de
%  l'orage (54) est « guguwar yamma », le vent d'ouest, et celui de
%  l'auberge au tableau 7 etait « iskar arewa ». C'est a mettre en
%  regard du tableau 3 : cette langue a emprunte TOUT son calendrier —
%  les sept jours a l'arabe, les douze mois a l'anglais — et pas une
%  seule de ses quatre directions. Le persan, au tableau 16, partage
%  sa rose des vents entre deux langues, شمال et جنوب arabes contre
%  خاور et باختر iraniens ; le haoussa n'a rien a partager. On
%  emprunte le temps, qui s'ecrit ; on gerde l'espace, qui se marche.
%
%  ET L'ARC-EN-CIEL (52) EST « bakan gizo », L'ARC DE L'ARAIGNEE.
%  Gizo est le heros fripon des contes haoussa, celui qui ruse et
%  gagne, et le ciel lui prete son arc apres l'orage. Aucune autre
%  colonne du releve ne nomme un phenomene par un personnage.
%
%  LE GRAIN DE LA PLANCHE EST DU BLE, ET LE BLE EST LE SEUL GRAIN QUE
%  LE HAOUSSA NOMME EN ARABE. gero le mil, dawa le sorgho, maiwa,
%  acca : tout ce qui pousse ici porte un nom d'ici. « alkama » le
%  ble est arabe, venu par la caravane avec la chose ; et le seigle
%  (43) n'a pas de nom du tout. C'est exactement la loi du tableau 3,
%  verifiee sur une troisieme matiere : ce qui a pousse ici se nomme
%  ici, ce qui est venu porte le nom d'ou il vient, ce qui n'est
%  jamais venu se decrit.
%
%  SEPTIEME REFUS DE LA COLONNE : LE NOYER (57). « goro » est la noix
%  de kola, et c'est la noix du monde haoussa : on l'offre pour
%  saluer, on la partage pour ouvrir une negociation, on l'envoie
%  pour demander une fille en mariage. Le mot serait tombe seul sous
%  la main devant un arbre a noix dont les feuilles jaunissent. Mais
%  la kola n'est pas la noix, et l'arbre de la planche n'est pas le
%  kolatier. On ecrit « itacen wolnat ».
%
%  MAIS LA BIERE EST « giya », SANS UN INSTANT D'HESITATION, ET C'EST
%  L'INVERSE. Le mot est celui du tableau 4, ou il rendait le vin de
%  la noce ; il designe en propre la biere de mil, brassee ici depuis
%  toujours, et ce tableau-ci parle justement de biere et de houblon.
%  La ou le persan avait a choisir, le haoussa avait le mot exact. Le
%  cidre suit : « giyar tuffa », la biere de pommes.
%
%  ET LA REGLE DE L'APOSTROPHE A SERVI, POUR LA PREMIERE FOIS. La
%  note (*) de ce tableau parle de la « ka’idar Ido », la regle
%  idiste, et le mot avait ete tape avec une apostrophe droite.
%  L'apostrophe haoussa est un COUP DE GLOTTE, une consonne pleine —
%  ka’ida, sa’a, jama’a, ma’aikata —, et l'orthographe de 2026 l'ecrit
%  ’ (U+2019) et non '. Douze regles avaient ete posees au tableau 1 ;
%  celle-ci attendait le huitieme tableau pour mordre une fois.
%
%  DEUX MOTS QUI REVIENNENT ET QUI COMPTENT. « ƙugiya » est le croc a
%  vetements du tableau 1 (58) et l'hamecon du pecheur ici : sixieme
%  paire d'homonymes de cette colonne, apres kaka, kwano, allo,
%  allura et keke. Et le gobelet (13) — dans le bloc aveugle, justement
%  — est « ƙoƙo », la demi-calebasse, qui nommait le crane au
%  tableau 2 et le de a coudre au tableau 6. Le haoussa nomme par la
%  FORME, et il ne se lasse pas.
% ===================================================================

%%K t08-noto-1 noto
\VUnotes{143.2mm}{%
\textit{(*) Domin wannan kalma ba ta da tabbaci : girbin lilin ne, ko
girbin inabi, ko girbin tuffa ? Wataƙila zai yi kyau a yi amfani da
«} \VUgras{mesono} \textit{» wadda aka gabatar a Mondo XIV, shafi na
242. Amma har yanzu Cibiyar ba ta karɓi wannan sabuwar saiwa ba,
tana kuma jiran tattaunawa bisa ga ka’idar Ido da aka saba.}%
}

%%K t08-apar-1 apar
\VUsaut{5.0mm}
\VUtitre{15.0pt}{60}{HOTO Na 8}
\VUsaut{3.0mm}
\VUcentre{13.2pt}{90}{\textit{Fage na Farko.}}
\VUsaut{3.0mm}
\VUpk{10.2pt}{40}{Girbi (*).}
\VUsaut{4.0mm}
\VUpk{10.2pt}{40}{Yanayin Ƙauye.}
\VUsaut{3.0mm}

%%K t08-c1-01-1 p
1. --- Tare da \VUgras{rani}, yanayin ƙauye ya sake canja. Irin da aka
danƙa wa layin noma ya tsira ; ƙaramin tsiro mai ganye ya fito daga
ƙasa ya kuma yi kan hatsi. Ƙwaya ta yi, sa’an nan ta nuna, har ma
yanzu ba a buƙatar komai sai girbi.

%%K t08-c1-02-1 p
2. --- \VUgras{Furannin ƙauye} sun buɗe. \VUgras{Furen shuɗi}
\textsuperscript{(1)}, da \VUgras{furen ja} \textsuperscript{(2)} da
\VUgras{furen ƙararrawa} \textsuperscript{(3)} suna haɗa wa launi ɗaya na
gonakin alkama bambancin \VUgras{furanninsu} sabbi masu farin ciki.

%%K t08-c1-03-1 p
3. --- Itatuwan ƴaƴan itace sun yi ado da ganye ; ƴaƴan itacen sun
nuna. Ƙaramin \VUgras{itacen ceri} \textsuperscript{(4)} yana cike da
kyawawan \VUgras{ceri} \textsuperscript{(5)} waɗanda yara suke tsinka
suna ci da jin daɗi ƙwarai.

%%K t08-c1-04-1 p
4. --- Yanzu dabbobi za su iya ci gaba da zama a sararin iska tsawon
yini.

%%K t08-c1-04-2 p
\VUgras{Ƴan tumaki} \textsuperscript{(6)} sun girma sun zama kyawawan
\VUgras{tumaki} \textsuperscript{(7)} da kyawawan \VUgras{raguna}
\textsuperscript{(8)} masu gashi mai kauri. Suna kiwon \VUgras{ciyawar}
\VUgras{makiyaya} \textsuperscript{(9)} cikin nutsuwa, makiyayar da
\VUgras{furannin fari} suka yi mata ado.

%%K t08-c1-04-3 p
Ba da daɗewa ba za a fara aske \VUgras{waɗannan} dabbobi domin a sami
\VUgras{gashinsu} wanda ake yin kaɗin tufafinmu da shi.

%%K t08-tit-2 sub
\VUsaut{5.0mm}
\VUpk{10.2pt}{40}{Makiyayi.}
\VUsaut{3.4mm}

%%K t08-c1-05-1 p
5. --- \VUgras{Makiyayi} \textsuperscript{(10)} ya yi kama da wanda bai
kula da su sosai ba. Hadarin da yake wucewa a sararin sama kaɗai yake
damunsa, \VUgras{mai kiwon tumaki} \textsuperscript{(13)} kuwa, wanda
yake kusanta \VUgras{salkarsa} \textsuperscript{(14)} da leɓɓansa, ya yi
kama da wanda ya fi shi rashin damuwa.

%%K t08-c1-06-1 p
6. --- Zafi yana matsawa ; domin \VUgras{kare} \textsuperscript{(15)}
ya zo, shi ma, ya kashe ƙishirwarsa ; yana lasar ruwan sanyi na
\VUgras{ƙaramin kogi} \textsuperscript{(16)} yana kuma tsoratar da
\VUgras{ƙaramin kwaɗo} \textsuperscript{(17)} mai tausayi wanda ya ɓoye a
cikin ciyawar gaɓa. Shi kuwa ya yi tsalle cikin ruwa mai tsabta inda
\VUgras{furannin ruwa} \textsuperscript{(18)} suke yin fure.

%%K t08-c1-07-1 p
7. --- \VUgras{Tsuntsun ruwa} \textsuperscript{(19)} yana yin wanka
kusa da \VUgras{ƙaramar marmaro} \textsuperscript{(20)} wadda take
ɓulɓula tsakanin duwatsu biyu, yana kuma ɓoye kansa a ƙarƙashin
\VUgras{ƙayayyuwa} \textsuperscript{(21)} da \VUgras{ƙayayyuwan fari}
\textsuperscript{(22)}.

%%K t08-tit-3 sub
\VUsaut{5.0mm}
\VUpk{10.2pt}{40}{Girbi.}
\VUsaut{3.4mm}

%%K t08-c1-08-1 p
8. --- Ga yaran ƙauye, girbin alkama lokaci ne mai daɗi ƙwarai. Suna
\VUgras{birgima}\textsuperscript{(24)} a kan \VUgras{tudun karan hatsi}
\textsuperscript{(23)} da farin ciki, suna taimakon \VUgras{masu girbi}
\textsuperscript{(26)}, kuma alhali waɗannan suna yanka \VUgras{gonar
alkama} \textsuperscript{(27)} da \VUgras{manyan laujensu}
\textsuperscript{(25)} da aka wasa da kyau a kan \VUgras{dutsen wasa}
\textsuperscript{(30)}, su suna tara alkamar suna kuma ɗaure ta a
\VUgras{dammuna} \textsuperscript{(31)} ko \VUgras{ƙananan dammuna}
\textsuperscript{(32)}.

%%K t08-c1-09-1 p
9. --- Wani lokaci akan yarda wa waɗanda suka fi girma a cikinsu su
yanka da \VUgras{lauje} \textsuperscript{(12)} \VUgras{kututturan
zinare} \textsuperscript{(28)} waɗanda suke ɗauke da \VUgras{kan hatsi}
\textsuperscript{(29)} masu nauyi cike da ƙwaya ; ƙanana kuwa, suna
lura da \VUgras{dabbobi} \textsuperscript{(33)} da suke kiwo a
\VUgras{makiyaya} \textsuperscript{(34)} a ƙarƙashin inuwar babbar
\VUgras{bishiya mai inuwa} \textsuperscript{(35)}, suna kama kyawawan
\VUgras{malam buɗe littafi} da \VUgras{tarunsu} \textsuperscript{(36)}.

%%K t08-c1-09-2 p
Har waɗansu suna hawa \VUgras{karusa} \textsuperscript{(37)} domin su
shirya dammunan da ake miƙa musu da \VUgras{sandar mai yatsa}
\textsuperscript{(38)}.

%%K t08-c1-10-1 p
10. --- A dukan \VUgras{fili} \textsuperscript{(39)} inda
\VUgras{tsuntsayen sama} \textsuperscript{(41)} suke tashi, ana cikin
hargitsi. Kowa yana aikin girbi. Amma, alhali matalauta suna ci gaba
da amfani da tsofaffin kayan aiki, wato \VUgras{lauje, manyan lauje}
da \VUgras{sandunan dako} \textsuperscript{(40)}, masu gonaki masu
arziki suna sa a yanka alkamarsu ko \VUgras{hatsin rai}
\textsuperscript{(43)} da \VUgras{injin yanka} \textsuperscript{(42)}.
Sa’an nan, tun da ba a buƙatar kai dammuna zuwa rumfar hatsi, akan yi
manyan \VUgras{tuduna na dammuna} \textsuperscript{(44)}.

%%K t08-c1-11-1 p
11. --- Sa’an nan akan kawo \VUgras{injin sussuka}
\textsuperscript{(47)} wanda \VUgras{injin tuƙi} \textsuperscript{(45)}
yake motsawa da \VUgras{igiyar injin} \textsuperscript{(46)}, haka nan
kuwa ake raba ƙwaya da \VUgras{karan hatsi,} ba tare da an bar gonar da
aka shuka ta ba.

%%K t08-tit-4 sub
\VUsaut{5.0mm}
\VUpk{10.2pt}{40}{Hadari.}
\VUsaut{3.4mm}

%%K t08-c1-12-1 p
12. --- Sau da yawa \VUgras{manyan hadari} \textsuperscript{(53)} sukan
tashi a lokacin girbi. Da farko ana jin cida na \VUgras{tsawa} daga
nesa ; \VUgras{guguwar yamma} \textsuperscript{(54)} takan fara busawa
tana kuma tanƙwara manyan bishiyoyin \VUgras{daji}
\textsuperscript{(49)} tana huci. \VUgras{Baƙaƙen gajimare}
\textsuperscript{(56)} suna kai wa sararin sama hari ; \VUgras{walƙiya}
\textsuperscript{(55)} mai ƙyalli tana ratsa su da karkatattun layi.
\VUgras{Ruwan sama} da wani lokaci \VUgras{ƙanƙara} sukan fara sauka,
suna murƙushe amfanin gonar da yake a shirye a yanka.

%%K t08-c1-13-1 p
13. --- Sau da yawa walƙiya takan ƙona gini. Haka nan bara rufin
\VUgras{gidan niƙar iska} \textsuperscript{(50)} ya zama toka, biyu
kuma daga cikin \VUgras{manyan fikafikansa} \textsuperscript{(51)} suka
farfashe.

%%K t08-c1-14-1 p
14. --- Bayan sa’o’i kaɗan, natsuwa takan dawo. \VUgras{Bakan gizo}
\textsuperscript{(52)} mai ƙyalli yakan bayyana a sararin sama, halitta
kuma takan sake ɗaukar rayuwarta da aka katse na ɗan lokaci. Mutanen
ƙauye, waɗanda suka gudu zuwa \VUgras{bukkoki} \textsuperscript{(48)},
sukan sake yin aiki suna kuma ƙoƙarin gyara barnar da suka sha da
ƙarfin hali.

%%K t08-tit-5 sub
\VUsaut{6.0mm}
\VUcentre{13.2pt}{90}{\textit{Fage na Biyu.}}
\VUsaut{3.6mm}

%%K t08-tit-6 sub
\VUpk{10.2pt}{40}{Farauta. --- Girbin inabi.}
\VUsaut{1.4mm}
\VUpk{10.2pt}{40}{Kamun kifi.}
\VUsaut{4.0mm}

%%K t08-c2-01-1 p
1. --- A kusan ƙarshen \VUgras{Agusta} ko tsakiyar \VUgras{Satumba},
bisa ga yankuna, ana buɗe farauta. Da ƙyar hasken rana na farko ya sa
\VUgras{ƙadangaru} \textsuperscript{(2)} su fito daga ramukansu, sai
\VUgras{mafarauci} \textsuperscript{(5)} ya kama hanya tare da
\VUgras{karnukansa} \textsuperscript{(4)}.

%%K t08-c2-02-1 p
2. --- Ya sa \VUgras{kayan farautarsa} \textsuperscript{(9)} masu ƙarfi
na kyalle mai kauri, ya sa \VUgras{murfin ƙafa} na fata a ƙafafunsa,
waɗanda za su sa shi ya iya tafiya cikin ciyawa mai ɗanshi inda
\VUgras{kwaɗon ƙasa} \textsuperscript{(1)} suke ɓoyewa ; ya ɗauki
\VUgras{bindigarsa} \textsuperscript{(6)} da \VUgras{jakar farautarsa}
\textsuperscript{(10)}, ga shi kuwa ya tafi.

%%K t08-c2-03-1 p
3. --- Himmar korar ta kai shi tsakiyar gonar inabi inda girbin inabi
yake da ƙarfi ƙwarai. Yaran manomin inabi, waɗanda a kullum suke lura
da awaki a kan hanya, yau sun bar su su yi kiwo yadda suka ga dama,
kuma bayan sun ɗaure tsohon \VUgras{bunsuru} \textsuperscript{(3)} a
gungume, shi wanda ba zai kasa amfani da ƴancinsa domin ya gudu ba, su
ma sun ba wa kansu rabo cikin jin daɗin. Ɗaya ya kama gungun inabi
daga \VUgras{kwandon girbi} \textsuperscript{(12)} yana kuma jin daɗin
zuba \VUgras{ruwansa} \textsuperscript{(14)} cikin \VUgras{ƙoƙo}
\textsuperscript{(13)} domin ya yi ruwan inabi da bai tuya ba. Ɗayan
kuwa ya sa \VUgras{kambin ganye} \textsuperscript{(15)} da ya saƙa yanzu
a kansa.

Kusa da su, \VUgras{tsuntsayen gida} \textsuperscript{(16)} marasa kunya
suna zuwa har gaban injin matsa inabi domin su saci inabin.

%%K t08-c2-04-1 p
4. --- Alhali yaran suna nishaɗi, maza suna aiki. \VUgras{Masu girbin
inabi} \textsuperscript{(18)} suna kai inabin cikin ɗakin ƙasa da
\VUgras{kwanduna na baya} \textsuperscript{(17)} ; \VUgras{mai yin
ganga} \textsuperscript{(19)} yana gyara \VUgras{gangunan giya}
\textsuperscript{(20)}, yana duba ko \VUgras{katakansu}
\textsuperscript{(22)} da \VUgras{ƙasansu} \textsuperscript{(21)} suna
cikin kyakkyawan hali, yana kuma musanya \VUgras{zoban ƙarfe}
\textsuperscript{(23)} da suka karye. Shi ne kuma ya yi aikin
\VUgras{manyan kwanuka} \textsuperscript{(25)} da ake zuba \VUgras{inabi}
\textsuperscript{(26)} a cikinsu domin su tuya.

%%K t08-c2-05-1 p
5. --- Idan tuyar ta ƙare, akan ɗiba giyar akan kuma ajiye saura a kan
\VUgras{injin matsa inabi} \textsuperscript{(28)}, ana matsawa a hankali
da \VUgras{babban sukuru} \textsuperscript{(29)} akan matse ruwan da
yake gudana cikin \VUgras{ƙaramin kwano} \textsuperscript{(32)}, har
yanzu kuma ana samun \VUgras{giyar inabi} \textsuperscript{(31)}.

%%K t08-tit-8 sub
\VUsaut{6.0mm}
\VUpk{10.2pt}{40}{Giya da giyar tuffa.}
\VUsaut{3.4mm}

%%K t08-c2-06-1 p
6. --- A bangon \VUgras{ɗakin ƙasa} \textsuperscript{(27)} akwai
reshen \VUgras{ganyen giya} \textsuperscript{(30)} yana hawa. A can tsiro
ne na ado kaɗai, amma a waɗansu ƙasashe, inda inabi yake da wuya
ƙwarai, ana noma ganyen giya domin a yi \VUgras{giya.}

%%K t08-c2-07-1 p
7. --- Ƙaramin \VUgras{itacen tuffa} \textsuperscript{(33)} yana cike
da tuffa waɗanda za su ba da \VUgras{giyar tuffa} mai kyau ƙwarai idan
sun nuna. A cikin \VUgras{kejinsu} \textsuperscript{(34)},
\VUgras{kanari} \textsuperscript{(35)} da \VUgras{tsuntsu mai launuka}
\textsuperscript{(36)} suna kallon tsuntsayen gida da suke yawo cikin
ƴanci da idon hassada.

%%K t08-c2-08-1 p
8. --- Har inda ido ya kai, a gefen tuddai, an jera \VUgras{sandunan
inabi} \textsuperscript{(40)} waɗanda suke ɗauke da kyawawan
\VUgras{kututturan inabi} \textsuperscript{(39)}, waɗanda
\VUgras{gungun inabi} masu nauyi suka yi wa ado.

%%K t08-c2-08-2 p
Tsawon shekara duka, \VUgras{manomin inabi} \textsuperscript{(42)} ya
yi aiki domin ya kula da \VUgras{gonar inabinsa} \textsuperscript{(38)},
yanzu kuwa an saka masa wahalarsa da girbi mai kyau. \VUgras{Mata masu
girbin inabi} \textsuperscript{(41)} suna cika kwanukan da aka ajiye a
kan karusar da \VUgras{alfadarai} \textsuperscript{(43)} suke ja, kowa
kuwa yana cikin farin ciki.

%%K t08-tit-9 sub
\VUsaut{5.0mm}
\VUpk{10.2pt}{40}{Kamun kifi.}
\VUsaut{3.4mm}

%%K t08-c2-09-1 p
9. --- Haka nan yake ga \VUgras{masu kamun kifi da ƙugiya}
\textsuperscript{(44)} waɗanda tun da safe suka zo suka zauna a gaɓar
ruwa da \VUgras{sandunan kamun kifinsu} \textsuperscript{(45)}.

%%K t08-c2-09-2 p
\VUgras{Kifi} \textsuperscript{(47)} yakan kama \VUgras{ƙugiya} nan take
da suka sa \VUgras{zarensu} \textsuperscript{(46)} cikin ruwa, kuma da
ƙyar suka sami lokacin sabunta \VUgras{koto,} sai wani sabon wanda aka
kama ya yi rawa a ƙarshen zaren.

%%K t08-c2-09-3 p
Duk da haka \VUgras{mai taru} \textsuperscript{(48)}, wanda daga
\VUgras{jirgin ruwansa} \textsuperscript{(49)} yake jefa \VUgras{taru}
zuwa tsakiyar \VUgras{kogi} \textsuperscript{(50)}, yakan kama kifi mafi
yawa.

%%K t08-c2-10-1 p
10. --- Kaka ita ce lokacin yawon shaƙatawa mai kyau : da ƙafa, da
\VUgras{hawan doki} \textsuperscript{(52)}, da \VUgras{hawan keke}
\textsuperscript{(53)}, da \VUgras{mota} \textsuperscript{(54)}.
\VUgras{Hanyoyi} \textsuperscript{(51)} suna da tsabta ana kuma amfani
da kwanaki masu kyau na ƙarshe, domin ga shi tuni \VUgras{alallaka}
\textsuperscript{(56)} suna taruwa a kan rufin gidaje domin su tashi da
dukan fikafikansu zuwa ƙasashe masu ɗumi. Ganyayyakin tsohon
\VUgras{itacen wolnat} \textsuperscript{(57)} suna yin rawaya,
\VUgras{wolnat} suna faɗuwa, dogayen \VUgras{bishiyoyi} kuwa da aka
tsara a \VUgras{gungu} \textsuperscript{(58)} a kan \VUgras{tudu}
\textsuperscript{(59)} ba da daɗewa ba za su rasa ganyensu.

%%K t08-c2-11-1 p
11. --- \VUgras{Rana} \textsuperscript{(60)} tana fitowa a makare,
kwanaki suna gajarta, ɗan sanyi mai kaifi yana fara jin sa da
\VUgras{safiya,} ga shi kuwa tuni garken \VUgras{tsuntsayen ƙaura}
\textsuperscript{(61)} suna barin ƙasashenmu marasa karɓar baƙi.

\VUsaut{6.0mm}
\VUfilet{22mm}
